\section{Results}
\label{sec:results}
%======================================================================
% CROSS-SCALE RESULTS: 2v2 and 6v6 are reported as ONE argument, organised by
% question, not as two standalone scale reports.
%
% 4v4 is deliberately EXCLUDED from this section (PI direction 2026-09-11).
%
% PENDING (one slot only): the 6v6 sealed crossover Delta_A / Delta_B in
% Sec.~\ref{sec:results-scale}. Search this file for "PENDING-6V6" -- it is the
% only number not yet sealed. Fill from
%   artifacts/strategic_demand/sppo/RUNG1_6V6_CROSSOVER_EVAL_RESULT.json
% (or ..._INTEGRITY_REQUIRED.json if a tie/reversal is flagged).
%======================================================================

\subsection{Evaluation protocol}
\label{sec:results-protocol}

Every number in this section comes from a \emph{sealed} evaluation artifact:
a seed block fixed before any episode ran, a terminal checkpoint chosen by a
budget rule rather than by score, and a one-shot write. No checkpoint, seed
block, severity level, or gate threshold in this section was selected after
seeing an outcome it could influence.

\paragraph{One criterion at both scales.}
Two strategy codes $z_0,z_1$ are \emph{payoff-specialized} only if each beats the
other against its intended opponent regime:
\[
\Delta_A=V(z_0,A)-V(z_1,A),
\qquad
\Delta_B=V(z_1,B)-V(z_0,B),
\]
with both point estimates strictly positive \emph{and} both
$\mathrm{LCB}_{95}(\Delta)>0$ under a paired percentile bootstrap over
evaluation seeds ($n_{\mathrm{boot}}=20{,}000$, $\alpha=0.05$, seed as the
resampling unit). We call this the \emph{specialization criterion}. It is applied
verbatim at $2$v$2$ and $6$v$6$; no scale receives a relaxed gate, a larger
$\alpha$, or a re-tuned bootstrap.

\paragraph{Why the two scales are directly comparable.}
The deployed actor is weight-shared across agents: team size changes the
observation batch dimension, not any parameter shape. The compression under test
is therefore \emph{numerically identical} at both scales---two independent
specialist actors occupy $6{,}947{,}184$ unique parameters, and tying only the
perception encoder reduces this to $3{,}612{,}784$, a $48.0\%$ reduction, at
$2$v$2$ and at $6$v$6$ alike (Table~\ref{tab:crossscale}). Scale changes the
coordination problem, not the parameter budget being saved. This isolates the
scientific question: does the same compression preserve the same complementary
specialization when six agents must coordinate instead of two?

\paragraph{Terminology used in figures and captions.}
These phrases are pre-specified, not informal judgments about whether a run
``looked good.''
\begin{itemize}
  \item \textbf{Satisfies / fails criterion} (figures may label these
  \textsc{pass}/\textsc{fail}): the criterion above is met, or is not. Point
  estimates alone never decide the label.
  \item \textbf{$\mathrm{LCB}_{95}$ / $\mathrm{UCB}_{95}$}: endpoints of a
  percentile bootstrap $95\%$ interval over matched evaluation seeds.
  \item \textbf{Detectable loss}: relative to the Share-0 expert-dispatch
  reference, a paired change
  $D_p=\Delta_p^{(\mathrm{cond})}-\Delta_p^{(\mathrm{Share\textrm{-}0})}$ is a
  detectable specialization loss on regime $p$ only if
  $\mathrm{UCB}_{95}(D_p)<0$ (the whole interval lies below zero).
  \item \textbf{Contested matchup}: neither side of a policy--opponent pairing is
  saturated, so specialization contrasts are informative rather than
  floor/ceiling artifacts.
  \item \textbf{Certified} (opponent-regime demand only): a preregistered
  strategic-demand protocol was satisfied under its sealed evaluation. It is a
  statement about the \emph{environment}, and does not by itself imply that any
  learned controller specializes.
\end{itemize}

\begin{table}[t]
\centering
\caption{Sealed evidence blocks. The same criterion, bootstrap, and sealing
discipline apply at both scales; the blocks differ in what each scale is asked to
establish. $2$v$2$ localizes \emph{where} sharing breaks specialization; $6$v$6$
tests whether the rung that survived localization also survives scale.}
\label{tab:crossscale}
\begin{tabular}{@{}lll@{}}
\toprule
\textbf{Evidence} & \textbf{$2$v$2$} & \textbf{$6$v$6$} \\
\midrule
Opponent-regime demand   & $n{=}192$, certified & $n{=}64$, certified \\
Independent experts      & $n{=}64$, criterion met & teachers only \\
Expert-dispatch (Share-0)& $n{=}128$, reference & not rebuilt \\
Shared encoder           & $n{=}128$, criterion met & $n{=}128$, sealed \\
Deeper sharing rungs     & $n{=}128$, erosion mapped & not ported \\
Unique actor parameters  & $6{,}947{,}184\!\to\!3{,}612{,}784$
                         & $6{,}947{,}184\!\to\!3{,}612{,}784$ \\
Robustness dose--response& $n{=}128$, 3 families & block reserved \\
\bottomrule
\end{tabular}
\end{table}

\begin{figure*}[t]
\centering
\includegraphics[width=\textwidth]{fig_cross_scale_2v2_6v6.pdf}
\caption{The cross-scale argument in three steps, at $2$v$2$ and $6$v$6$ under one
criterion. (a)~Opponent-regime demand is certified at both scales and does not
weaken with team size---the Regime-$B$ contrast is larger at $6$v$6$.
(b)~Matched-state distillation into a single shared-encoder controller retains
$\approx99\%$ of the teachers' branch divergence at both scales, under a
numerically identical $48.0\%$ reduction in unique actor parameters.
(c)~Whether that compressed policy still satisfies the specialization criterion:
sealed and satisfied at $2$v$2$; the $6$v$6$ cell is drawn as pending because its
evaluation is not yet sealed, and is never estimated.
Error bars are $95\%$ bootstrap CIs.}
\label{fig:crossscale}
\end{figure*}

%----------------------------------------------------------------------
\subsection{The opponent regimes demand different strategies at both scales}
\label{sec:results-demand}

Before asking whether a controller \emph{discovers} two strategies, we establish
that the environment \emph{requires} them. A preregistered opponent-conditioned
protocol scores scripted guard-style and breach-style team responses against each
regime on paired seeds, and certifies demand only if each response strictly beats
the other on its own regime with $\mathrm{LCB}_{95}>0$.

At $2$v$2$ ($n{=}192$ paired seeds), the payoff gaps are
\[
\Delta_G(A)=+0.297\;[+0.203,+0.391],
\qquad
\Delta_B(B)=+0.453\;[+0.385,+0.526],
\]
with concealment gates satisfied on both regimes ($p_C(A)=0.943$,
$p_C(B)=0.917$), so the regime is not identifiable from a single observation.

At $6$v$6$ ($n{=}64$ paired seeds), the poles are size-normalized---the defensive
trigger scales with team size, and the guard response distributes
$\lceil N/2 \rceil$ defenders rather than stacking them on one intruder---and the
same protocol certifies demand with
\[
\Delta_G(A)=+0.266\;[+0.094,+0.438],
\qquad
\Delta_B(B)=+0.625\;[+0.484,+0.750].
\]
Two features matter for the scaling argument (Fig.~\ref{fig:crossscale}a). First,
demand does not weaken with team size: the Regime-$B$ contrast is substantially
\emph{larger} at $6$v$6$ ($+0.625$ vs.\ $+0.453$), driven by guard-style play
collapsing against the breach regime ($0.078$ win rate) while breach-style play
holds ($0.703$). Second, both regimes remain contested rather than saturated at
both scales, so the downstream specialization contrasts are informative.

Certified demand is a property of the environment. It does not predict that a
learned controller will discover the two strategies, and the remainder of this
section is about that gap.

%----------------------------------------------------------------------
\subsection{Independent experts realize the demand as complementary payoff}
\label{sec:results-experts}

Independent PPO experts $\pi_A,\pi_B$ trained against the two regimes satisfy the
specialization criterion at $2$v$2$ on a held-out $64$-seed block:
\[
\Delta_A = +0.172\;[+0.016,+0.328],
\qquad
\Delta_B = +0.297\;[+0.125,+0.469],
\]
while a mixed-opponent generalist trained on the same budget attains
$V(\pi_G,A)=0.375\;[0.265,0.500]$ and $V(\pi_G,B)=0.500\;[0.375,0.625]$---it
does not recover either specialist's own-regime advantage
(Fig.~\ref{fig:winrate-context}).

To separate ``the strategies exist'' from ``the evaluation path can express
them,'' we score a Share-0 reference: bit-exact $z$ dispatch between the verbatim
experts, with zero tied modules. On the matched $128$-seed block it satisfies the
criterion with
\[
\Delta_A = +0.289\;[+0.164,+0.406],
\qquad
\Delta_B = +0.258\;[+0.141,+0.375],
\]
establishing the ceiling that any shared policy is measured against.

At $6$v$6$, specialists are trained under the identical recipe, budget, and
checkpoint rule and serve as distillation teachers. We deliberately do not
re-certify expert crossover or rebuild the Share-0 ladder at $6$v$6$: the
mechanism and its breaking point were localized at $2$v$2$, and the $6$v$6$ track
is a scaling test of the rung that survived, not a rediscovery of the ladder.
This is the one structural asymmetry between the scales, and it is a scope
decision rather than a missing measurement.

%----------------------------------------------------------------------
\subsection{Distillation fidelity licenses the payoff question---it does not
answer it}
\label{sec:results-fidelity}

Matched-state distillation transfers the two experts into discrete modes $z_0,z_1$
of a single controller whose perception encoder is tied ($96.0\%$ of one branch's
parameters) and whose remaining capacity stays private per mode.

Imitation fidelity is high at both scales. At $2$v$2$, holdout argmax agreement
with the source experts is $0.967/0.989$, and the student's branch
Jensen--Shannon divergence is $0.178$ against the experts' own $0.179$. At
$6$v$6$, holdout agreement is $0.963/0.961$ with holdout
$\mathrm{KL}=0.054/0.094$, and the student reproduces $0.478$ of the teachers'
$0.484$ branch divergence. In both cases the student retains roughly $99\%$ of the
behavioral separation present in its teachers (Fig.~\ref{fig:crossscale}b), and a
save/load round trip reproduces logits exactly. Note that the $6$v$6$ teachers are
far more behaviorally distinct in absolute terms ($0.484$ vs.\ $0.179$), so the
compressed policy is asked to preserve a substantially wider separation at the
larger scale.

These are necessary conditions, not the claim. The $2$v$2$ ladder supplies the
counterexample that forces this distinction: Share-Macro retains $0.936/0.969$
agreement---high by any imitation standard---and still \emph{fails} the
specialization criterion (Fig.~\ref{fig:hierarchy}). Action-level agreement and
nonzero action-distribution separation can both survive while closed-loop payoff
specialization does not. Fidelity is therefore reported as a precondition for
asking the payoff question at each scale, and the criterion is what answers it.

\begin{figure}[t]
\centering
\includegraphics[width=\columnwidth]{fig_2v2_measurement_hierarchy.pdf}
\caption{Measurement hierarchy: high imitation and nonzero action-distribution
separation do not imply closed-loop payoff specialization. Share-Macro retains
strong expert agreement yet fails the criterion via
$\mathrm{LCB}_{95}(\Delta_A)=0$.}
\label{fig:hierarchy}
\end{figure}

%----------------------------------------------------------------------
\subsection{Sharing only perception preserves specialization; sharing more
erodes it}
\label{sec:results-ladder}

At $2$v$2$ we sweep how much of the controller is tied, holding the distillation
recipe fixed, and score each rung on the same $128$ seeds
(Table~\ref{tab:sharing-results}, Fig.~\ref{fig:sharing-ladder}).

Sharing only the perception encoder preserves the criterion on both regimes with
no detectable paired loss against Share-0 ($D_A=-0.023$, $D_B=-0.094$; both
$\mathrm{UCB}_{95}\ge0$), while removing $48.0\%$ of the unique actor parameters.
Tying the backbone still satisfies the absolute criterion but now shows a
detectable Regime-$A$ loss, with on-regime performance falling and cross-regime
performance rising---the signature of homogenization rather than uniform
capacity loss. Tying the macro layer fails the criterion outright
($\mathrm{LCB}_{95}(\Delta_A)=0$) and additionally falls outside the imitation
fidelity band, so its degradation is reported as sharing-associated and
fit-confounded rather than uniquely attributed to homogenization.

The engineering consequence is sharp: essentially all of the parameter saving is
already captured by sharing perception alone. Backbone and macro sharing add only
about $1.5$ percentage points of further reduction while progressively destroying
the property the compression was meant to retain.

\begin{table}[t]
\centering
\caption{Parameter--specialization tradeoff on the matched $128$-seed set
($2$v$2$). Unique actor parameters are sealed construction counts. Reduction is
relative to Share-0. $D_A$/$D_B$ are within-seed differences vs.\ Share-0;
``detectable loss'' means $\mathrm{UCB}_{95}(D)<0$.}
\label{tab:sharing-results}
\begin{tabular}{@{}lrrlll@{}}
\toprule
\textbf{Condition} & \textbf{Params (M)} & \textbf{Red.} &
$\boldsymbol{\Delta_A}$ & $\boldsymbol{\Delta_B}$ &
\textbf{Criterion / $D$} \\
\midrule
Share-0
  & $6.947$ & $0\%$
  & $+0.289\;[+0.164,+0.406]$
  & $+0.258\;[+0.141,+0.375]$
  & PASS (reference) \\
Share-Encoder
  & $3.613$ & $\mathbf{48.0\%}$
  & $+0.266\;[+0.148,+0.383]$
  & $+0.164\;[+0.047,+0.281]$
  & PASS; no detectable loss \\
Share-Backbone
  & $3.509$ & $49.5\%$
  & $+0.148\;[+0.023,+0.266]$
  & $+0.141\;[+0.016,+0.266]$
  & PASS; detectable loss on $A$ \\
Share-Macro
  & $3.508$ & $49.5\%$
  & $+0.125\;[0.000,+0.250]$
  & $+0.156\;[+0.031,+0.281]$
  & FAIL ($\mathrm{LCB}_{95}(\Delta_A)=0$) \\
\bottomrule
\end{tabular}
\end{table}

\begin{figure*}[t]
\centering
\includegraphics[width=\textwidth]{fig_claim_a_2v2.pdf}
\caption{Causal chain at $2$v$2$.
(a)~Independent PPO experts establish complementary payoff strategies.
(b)~Matched-state distillation transfers them into discrete modes of one
shared-encoder controller ($z_0\!\sim\!\pi_A$, $z_1\!\sim\!\pi_B$).
(c)~Sealed $\Delta_A,\Delta_B$ with $95\%$ CIs across progressive sharing
($n{=}128$). Share-0 is a bit-exact expert-dispatch control; the remaining rungs
share one distillation recipe and differ only in which modules are tied.}
\label{fig:claim-a}
\end{figure*}

\begin{figure*}[t]
\centering
\includegraphics[width=\textwidth]{fig_sharing_ladder_2v2.pdf}
\caption{Exact sealed $\Delta_A,\Delta_B$ (left) and paired $D$ vs.\ Share-0
(right) with $95\%$ CIs ($n{=}128$). Share-Backbone still satisfies the absolute
criterion but already shows detectable Pole-$A$ loss; Share-Macro fails because
$\mathrm{LCB}_{95}(\Delta_A)=0$.}
\label{fig:sharing-ladder}
\end{figure*}

\begin{figure*}[t]
\centering
\includegraphics[width=\textwidth]{fig_absolute_winrate_context_2v2.pdf}
\caption{Absolute win rates for the mixed-opponent generalist $\pi_G$, the
independent experts, and the shared-encoder modes against both regimes. The
crossover pattern is visible when A-specialized policies win more on Regime~$A$
and B-specialized policies win more on Regime~$B$. Absolute context only;
specialization is decided by the paired $\Delta$ criterion. Error bars are $95\%$
bootstrap CIs ($\pi_G$/experts $n{=}64$; shared encoder $n{=}128$).}
\label{fig:winrate-context}
\end{figure*}

%----------------------------------------------------------------------
\subsection{Does the preserved rung survive scale?}
\label{sec:results-scale}

The $2$v$2$ ladder identifies exactly one configuration worth carrying forward:
shared perception, private strategy capacity. At $6$v$6$ we rebuild that single
configuration---same architecture, same distillation recipe, same $48.0\%$
parameter reduction---and evaluate it against the same criterion on a fresh
$128$-seed block, with $z_0$ and $z_1$ dispatched from one checkpoint
(Fig.~\ref{fig:crossscale}c).

% ---------------------------------------------------------------- PENDING-6V6
% The sealed 6v6 crossover evaluation is the only number in this section that is
% not yet available. Evaluation launched 2026-09-11T18:33Z on seeds
% 13640001..13640128 (512 episodes, ~62 s/ep). On completion, replace the
% placeholder below with the sealed values from
%   artifacts/strategic_demand/sppo/RUNG1_6V6_CROSSOVER_EVAL_RESULT.json
%     -> PRIMARY_GATE.delta_A {mean, lcb95, ucb95}
%     -> PRIMARY_GATE.delta_B {mean, lcb95, ucb95}
%     -> PRIMARY_GATE.passes
% If RUNG1_6V6_CROSSOVER_EVAL_INTEGRITY_REQUIRED.json is written instead, a tie or
% reversal was flagged and a row-level audit is required BEFORE any sentence here
% is written.
% ----------------------------------------------------------------------------
\noindent
\textbf{[PENDING-6V6 --- sealed crossover in progress; do not typeset.]}
\[
\Delta_A = \texttt{<pending>},
\qquad
\Delta_B = \texttt{<pending>}.
\]

\noindent
The interpretation is pre-committed and does not depend on the outcome. If the
criterion is satisfied, the same shared-perception compression that preserved
complementary specialization with two agents also preserves it with six, and the
$48\%$ deployment saving transfers to a materially harder coordination problem.
If it is not satisfied, the compression result is scale-bounded: the location of
the sharing boundary established at $2$v$2$ does not by itself license deployment
at larger team sizes, and the boundary must be re-established per scale. Either
outcome is reported against the same criterion, and no rung beyond shared
perception is evaluated at $6$v$6$ in this work.

%----------------------------------------------------------------------
\subsection{What the latent modes actually do}
\label{sec:results-behavior}

Three complementary views establish that the modes correspond to recognizable
team strategies rather than to arbitrary policy perturbations: matched
trajectories, role-allocation proxies, and the payoff contrasts already reported.

Under identical initial conditions with only $z$ changed, the team paths diverge
at tick~$4$ and reach different terminal outcomes, with terminal scores matched
against the sealed evaluation rows (Fig.~\ref{fig:traj}); matched frames at a
fixed tick show the same seed under both codes and both regimes
(Fig.~\ref{fig:qual}). Role-allocation telemetry
(Figs.~\ref{fig:roles}--\ref{fig:roles-ext}) is reported as an exploratory
characterization at $n{=}24$ per cell: coarse means are largely regime-driven,
with smaller within-regime gaps between codes. We therefore treat role proxies as
descriptive support for the trajectory evidence, never as a specialization gate.

\begin{figure*}[t]
\centering
\includegraphics[width=0.92\textwidth]{fig_trajectory_strip_2v2.pdf}
\caption{Matched seed $11960003$, shared-encoder policy: identical initial
conditions within each panel; only forced $z$ differs. Solid~$z_0$, dashed~$z_1$;
stars mark first action divergence ($t{=}4$). Terminal scores match the sealed
evaluation rows. Visual evidence that changing $z$ changes team behavior---not a
substitute for the statistical criterion.}
\label{fig:traj}
\end{figure*}

\begin{figure*}[t]
\centering
\includegraphics[width=0.95\textwidth]{fig_qualitative_latent_2v2.pdf}
\caption{Matched-frame grid at tick~$30$ for the same seed and shared-encoder
policy. Columns differ only in forced $z$; rows differ only in opponent regime.}
\label{fig:qual}
\end{figure*}

\begin{figure*}[t]
\centering
\includegraphics[width=\textwidth]{fig_role_allocation_2v2.pdf}
\caption{Role-allocation proxies under forced $z_0$ vs.\ $z_1$ on both regimes
(defender count, attacker count, attack/defense ratio). Episode-mean telemetry
from an exploratory characterization ($n{=}24$/cell; not a specialization gate).}
\label{fig:roles}
\end{figure*}

\begin{figure*}[t]
\centering
\includegraphics[width=\textwidth]{fig_latent_behavior_2v2.pdf}
\caption{Extended role telemetry (same exploratory source as
Fig.~\ref{fig:roles}): defenders, attackers, intercept near enemy carrier, carrier
escort, and attack/defense ratio under forced $z_0$ vs.\ $z_1$ ($n{=}24$/cell).
Regime effects dominate; within-regime gaps between codes remain modest.}
\label{fig:roles-ext}
\end{figure*}

%----------------------------------------------------------------------
\subsection{Robustness as a sim-to-real diagnostic}
\label{sec:results-robustness}

Specialization that exists only under noiseless sensing is not a deployable
property. We therefore re-evaluate the preserved $2$v$2$ configuration under three
deployment-relevant perturbation families at a preregistered severity ladder,
scoring the same criterion at each dose on the same $128$ seeds
(Fig.~\ref{fig:robust}). Severities were fixed in the protocol before any
robustness episode ran, and no severity was added after seeing a result.

The three families dissociate cleanly. Localization noise leaves the criterion
satisfied at every tested dose, with paired changes relative to nominal
statistically indistinguishable from zero---specialization is genuinely robust to
this channel across the tested range. Motion error fails the joint criterion at
\emph{every} nonzero dose, always through Regime~$B$, and---critically---the
failure does not grow with severity: the mildest dose already reproduces it, and
no paired degradation interval excludes zero. The honest description is a
persistently fragile Regime-$B$ margin under any motion perturbation, not a
dose-dependent breakdown. Control latency behaves differently again: it is the one
family with a monotone dose response, concentrated entirely in Regime~$A$, whose
margin falls steadily as latency increases and reaches the boundary of statistical
detectability at the largest authorized dose.

That the three families fail in three different ways---not at all, immediately and
flatly, and monotonically---is itself the finding. It means ``robustness'' is not
a single scalar property of the compressed controller, and that a deployment
decision must be made per sensing channel rather than from an aggregate score.

\begin{figure*}[t]
\centering
\includegraphics[width=\textwidth]{fig_robustness_delta_dose.pdf}
\caption{Dose--response of specialization under deployment-relevant perturbations
($n{=}128$): localization (GPS-like), motion (currents/actuators), and control
latency (communications). \textsc{pass}/\textsc{fail} mark the specialization
criterion at each severity. Localization remains specialized at every tested dose;
motion breaks Regime-$B$ specialization from the first nonzero dose; latency shows
monotone Regime-$A$ erosion.}
\label{fig:robust}
\end{figure*}

%----------------------------------------------------------------------
\subsection{Scope}
\label{sec:results-scope}

Three limits bound these claims. First, the sharing ladder is localized at $2$v$2$
only; $6$v$6$ evaluates the single preserved rung, so the \emph{position} of the
sharing boundary is established at one scale and tested, not re-derived, at the
other. Second, robustness is characterized at $2$v$2$ across three families and a
preregistered severity ladder; the corresponding $6$v$6$ seed block is reserved
but not yet spent, so no cross-scale robustness claim is made. Third,
strategy-preserving fine-tuning interventions (formulated in the supplementary
material) did not recover the criterion under frozen matched protocols, and are
not credited for any result above.
